\documentclass[12pt, a4paper]{article}

% --- Paquetes ---
\usepackage[utf8]{inputenc}
\usepackage[spanish]{babel}
\usepackage{amsmath, amssymb}
\usepackage{geometry}
\usepackage{fancyhdr}
\usepackage{xcolor}
\usepackage{enumitem}
\usepackage{titlesec}

% --- Configuración de Página ---
\geometry{top=2.5cm, bottom=2.5cm, left=2.5cm, right=2.5cm}
\pagestyle{fancy}
\fancyhf{}
\lhead{Planificación del Póster Académico}
\rhead{Ciencia de Datos - Externado}
\cfoot{\thepage}

% Estilos de Título
\titleformat{\section}{\large\bfseries\color{darkgray}}{}{0em}{}[\titlerule]

% --- Título del Documento ---
\title{Estructura y Planificación para el Póster de Investigación \\ \large ``Integración de sentimiento de mercado en agentes de RL para la optimización de trading''}
\author{Julian Duarte \quad Julian Jimenez}
\date{Abril 2026}

\begin{document}

\maketitle

\section{Información General del Proyecto}
\begin{itemize}
    \item \textbf{Título}: Integración de sentimiento de mercado en agentes de RL para la optimización de trading.
    \item \textbf{Autores}: Julian Duarte, Julian Jimenez.
    \item \textbf{Tutores}: Prof. Julián Fernando Sánchez, Prof. Juan Carlos Urueña.
    \item \textbf{Semillero}: Semillero de Aprendizaje por Refuerzo.
    \item \textbf{Afiliación}: Universidad Externado de Colombia, Departamento de Matemáticas.
\end{itemize}

\section{Lineamientos Generales del Póster}

\begin{itemize}
    \item \textbf{Extensión Total:} Máximo entre 800 y 1,000 palabras para garantizar legibilidad y balance visual.
    \item \textbf{Tipografía (Mínimos):}
        \begin{itemize}
            \item Cuerpo de texto: $\ge$ 24 pt.
            \item Títulos de sección: $\ge$ 48 - 54 pt.
            \item Título principal: $\ge$ 80 - 120 pt.
            \item Estilo: Sans-serif (ej. Arial, Montserrat, Calibri).
        \end{itemize}
    \item \textbf{Lenguaje:} Académico, preciso y sobrio. Evitar metáforas innecesarias o lenguaje comercial.
    \item \textbf{Balance Visual:}
        \begin{itemize}
            \item 40-50\%: Gráficos, esquemas y diagramas.
            \item 20-30\%: Texto descriptivo.
            \item 25\%: Espacio en blanco (márgenes y separación entre bloques).
        \end{itemize}
    \item \textbf{Calidad de Imagen:} Resolución mínima de 300 dpi para impresión en gran formato.
\end{itemize}

\section{Contenido y Descripción por Secciones}

\subsection{1. Introducción y Contexto}
\textbf{Guía Científica:} Presentar el ``gancho''. Debe responder a ¿qué se está estudiando?, ¿por qué es importante? y ¿cuál es el vacío de conocimiento?
\begin{itemize}
    \item \textbf{Tono:} Conciso y motivador. Máximo 100-150 palabras.
    \item \textbf{Visual:} Imagen representativa de la volatilidad o psicología de mercado.
    \item \textbf{Contenido:} [Problema de los sesgos emocionales en el trading minorista de criptomonedas].
\end{itemize}

\subsection{2. Objetivo(s)}
\textbf{Guía Científica:} Declaración directa de lo que se espera lograr usando verbos en infinitivo.
\begin{itemize}
    \item \textbf{Tono:} Directo y puntual.
    \item \textbf{Visual:} Uso de viñetas claras.
    \item \textbf{Contenido:} [Desarrollar y validar teóricamente un agente de Q-Learning que integre el Fear \& Greed Index].
\end{itemize}

\subsection{3. Metodología (MDP)}
\textbf{Guía Científica:} Visión general del diseño y las herramientas (Binance API, Alternative.me API).
\begin{itemize}
    \item \textbf{Tono:} Técnico-pedagógico.
    \item \textbf{Visual:} **DIAGRAMAS TÉCNICOS:** Ciclo de RL, Vector de Estado discretizado y Función de Recompensa penalizada por riesgo.
    \item \textbf{Contenido:} [Formulación matemática del MDP y lógica de actualización del Q-Value].
\end{itemize}

\subsection{4. Resultados Esperados y Evaluación (Fase Teórica)}
\textbf{Guía Científica:} Dado que el proyecto está en fase de validación teórica, se deben proyectar los hallazgos y definir las métricas de éxito.
\begin{itemize}
    \item \textbf{Tono:} Prospectivo y riguroso. "Se espera que...", "El modelo será evaluado mediante...".
    \item \textbf{Visual:} Mock-ups de gráficas de rendimiento esperadas (ej. Sharpe Ratio estimado), tablas de métricas de riesgo (Drawdown) y comparativa hipotética vs. estrategia pasiva.
    \item \textbf{Contenido:} [Criterios de validación: Ratio de Sharpe, Máximo Drawdown y Resiliencia en periodos de pánico colectivo].
\end{itemize}

\subsection{5. Conclusiones y Aplicación}
\textbf{Guía Científica:} Interpretar los alcances teóricos y proponer proyecciones (fase de implementación real).
\begin{itemize}
    \item \textbf{Tono:} Reflexivo.
    \item \textbf{Visual:} Lista de "Puntos clave" del aporte de la investigación.
    \item \textbf{Contenido:} [Impacto de las finanzas conductuales en la automatización del trading].
\end{itemize}

\end{document}
